% UNINTEGRATED WORKING APPENDIX FRAGMENT, 24 October 2026.
% Same general compact paper family; no standalone manuscript or compilation.
% Complete endpoint transcription from the accepted whole-A18 source chain.
% Source preparation only: independent transcription review and integration
% are pending; the frozen foundation batch and canonical editions are unchanged.

\section{The whole-\(A_{18}\) endpoint argument}
\label{app:wholeA18}
Here the entire diagram at the original farthest vertex is
\(A_{18}\), with no spectators. Fix that vertex \(p\) and the
original strict interior point \(z\), and write
\[
 p^2=-1,\qquad z=E(p-x),\qquad E>0,\quad x\perp p,
 \quad \norm{x}^2<1.
\]
The spatial space is the zero-sum subspace of \(\mathbb R^{19}\).
Its unit incident normals and duals satisfy
\[
 n_i=(e_i-e_{i+1})/\sqrt2,\qquad
 \ip{\omega_i}{\omega_j}
   =\frac{2\min(i,j)(19-\max(i,j))}{19},\qquad
 D_0=\omega_1^2=\omega_{18}^2=36/19.
\]
Interiority gives \(x_1>\cdots>x_{19}\).
Let the actual endpoint opposites be the walls terminating the
edges obtained by releasing \(n_1\) and \(n_{18}\).
Compactness makes these endpoints finite. For each opposite \(c\),
its positive time is \(A=-\ip p c>0\).
We retain every original wall throughout. Reversing the coordinate
formulas gives their right-end versions; no chamber symmetry or
arithmeticity is assumed. Distinctness of the two opposites, when
their mutual native gap is used, will follow from their actual contacts.

\subsection{The reference under the full norm bound}
The \(N=19\) central inventory and helper nonadjacency
above precede any contrary energy assumption.
Its near-four roles can also be excluded under
only \(\norm x<1\), using their sharper actual mates.
For a left near-four at \(i=5,6\), release \(j=i-1\).
The positive replacement has diagram
\(B_{20-i}+A_{i-2}\) and cosine
\[
 C_j=\sqrt{2(i-1)/19}
     \bigl((20-i)/2+(19-i)/\sqrt2\bigr)>11.
\]
Its normalized height exceeds \(\sqrt{5/6}>.91\).
At the other end a mate or a strong \(13/14\)
has demand \(>209/270>.774\).
Their correlation is at most
\(\sqrt{15/91}<.41\).
Cauchy contradicts \(1.684^2=2.835856>2.82\).
Thus all central weak sites are nonadjacent pure
helpers already under the full norm bound.

Use the same \(N=19\) rational heights and chord
omissions. If \(k\) is the number of weak end
sites \(5,14\), one reference \(w\) has energy
\[
 E_0\geq.9240352,\qquad E_1>.8937,\qquad E_2>.86414.
\]
Its left and right block lengths \(m\) are \(5\)
or \(6\), with signed means
\[
 b_5=H_5/(5\sqrt2)\in(.29,.298),\qquad
 b_6=H_6/(6\sqrt2)\in(.25,.266).
\]
The separate boundary dual prices satisfy
\[
 a_5\geq(.4208-.151)/2=.1349,\qquad
 a_6\geq(451/1200-.108)/2>.1339.
\]
Every omission is included and can only increase
the adjacent available price.

For an actual endpoint pressure \(\ip{x}{V}\geq K\),
subtract \(B=\sum b_i\omega_i\), with \(b_i\geq0\) supported
on the retained boundary duals, so that \(Y=V-B\) has zero mean
on each end block. Write \(w=\sum a_i\omega_i\), with \(a_i\geq0\).
Every retained height satisfies \(\ip x{\omega_i}\geq H_i
=\ip w{\omega_i}\), and \(\ip wY=0\).
For \(\mu\geq0\), the individual conditions \(\mu b_i\leq a_i\)
give
\[
 \ip x{w+\mu Y}
 =\sum_i(a_i-\mu b_i)\ip x{\omega_i}+\mu\ip xV
 \geq w^2+\mu(K-\ip wB).
\]
Applying \(\norm{x-(w+\mu Y)}^2\geq0\) therefore gives the
reference energy plus the gain
\[
 2\mu(K-\ip{w}{B})-\mu^2Y^2.
\]
Opposite-block residuals have zero cross.
For a coordinate pressure \(x_{\rm end}\geq t\),
the legally optimized reference gain is
\(m(t-b_m)^2/(m-1)\).
For all thresholds used below, the debit is
less than \(.061\), below each separate price.
If either coordinate were \(\geq.603\), the
three energy bounds plus this gain would exceed
\(1\); the tight case is
\[
 .86414+(6/5)(.603-.266)^2=1.0004228.
\]
Therefore \(x_1,-x_{19}<.603\).
The genuine endpoint-edge inequality and
\(.603^2<90/247=(18/19)(5/13)\) imply
\begin{equation}\label{app:eq:A18shortends}
 C_L,C_R<9/4.
\end{equation}
This bounds actual edge cosines, not cap owns.

\subsection{Complete local isolated-endpoint gains}
An isolated left endpoint has
\(c=A p-g\omega_1,\ A^2=D_0g^2-1>0\).
Own zero, \(3,4\) has nonpositive time; own
\(6\) and parallel have the positive actual
nulls \(\sqrt3c+2n_1+n_2\) and \(c+n_1\).
The full remaining list is finite \(5\), finite
\(7\), finite at least \(8\), or dotted.
For a pressure normalized to raw coordinate
sum \(\sqrt2\), subtract \(\omega_m/m\);
the residual is zero-sum on the relevant block.
Every price below is paid before applying its
actual pressure. The resulting complete gain
table is
\[
\begin{array}{c|c|c|c}
\text{own}&\mu&\text{gain on }m=5>&
 \text{gain on }m=6>\\ \hline
5&.6&.02904&.059\\
7&.4&.05152&.07677\\
\text{finite }\geq8&1/4&.0471&157/2400\\
\text{dotted}&1/4&.06&4699/60000.
\end{array}
\]
We prove every row, including continuous dotted
owns and the unbounded finite-label tail.

For dotted \(g>1\), put \(u=g^{-2}\in(0,1)\).
The actual positive core \(c+n_1/g\) gives
\[
 V=\omega_1-un_1,\quad
 K(u)=\sqrt{D_0-u}-\sqrt{u(1-u)},\quad
 Y^2=2-2/m-2u+u^2.
\]
It has square \(u-1<0\), positive time,
zero cap product and negative old-\(1\)
product; positivity checks every unknown wall.
At \(\mu=1/4\) the gain is
\[
 G_m(u)=\tfrac12(K(u)-H_m/m)
        -(2-2/m-2u+u^2)/16.
\]
Its shift is exactly
\(G_6-G_5=1099/60000\).
To prove \(G_5>.06\) uniformly, it suffices that
\[
 K(u)>.7408-u/4+u^2/8.
\]
The upper bound
\(\sqrt{u(1-u)}\leq1/2-(u-1/2)^2\)
follows by positive squaring, with residual
\((u-1/2)^4\).
Use lower radical values \(1.376,1.1809,1.0699,.9459\)
at \(u=0,1/2,3/4,1\); their squares are
\(1.893376,1.39452481,1.14468601,.89472681\),
strictly below \(D_0-u\).
Concavity places \(\sqrt{D_0-u}\) above the
successive chords. The three resulting
strict quadratic lower bounds are
\[
\begin{array}{c|l}
[0,1/2]&.3852-1.1402u+(7/8)u^2\\
[1/2,3/4]&
(7/8)(u-597/875)^2+8357/1750000\\
[3/4,1]&.4511-1.246u+(7/8)u^2.
\end{array}
\]
The first decreases to \(677/20000>0\);
the last increases from \(703/80000>0\).
This proves the entire dotted row.

For finite \(g<1\), put \(u=g^{-2}\), \(d=D_0-1=17/19\).
The actual finite projection onto \(c,n_1\)
gives
\[
 V=\omega_1-n_1,\quad
 K_f=\frac{d-\sqrt{d(u-1)}}{\sqrt{D_0-u}},\quad
 Y^2=1-2/m.
\]
Writing \(s=\sqrt{u-1}\),
\(K_f^2=d(\sqrt d-s)/(\sqrt d+s)\), increasing
with \(g\).
For every native own at least \(8\),
\(s\leq\sqrt2-1<.414214\).
Using \(.8947<d<.8948,\sqrt d>.9459\),
\[
 (d-.59^2)\sqrt d>.51702894,\qquad
 (d+.59^2)s<.5148265806.
\]
Hence \(K_f>.59\). Quarter-payment gives
\(G_{f,5}>.0471\) and
\(G_{f,6}>.0471+1099/60000=157/2400\).

For native \(5\), let \(u=6-2\sqrt5\) and
\(b=D_0-3/2=15/38\).
Its actual \(H_4\) correction
\(\eta=c+g(3n_1/2+n_2+n_3/2)\)
has square \(\delta=1-3g^2/2>0\).
Use the actual finite point
\(p+(A/\delta)\eta\), not \(\eta\).
The pressure is
\[
 V=\omega_4/4,\qquad
 K_5=\frac{b-\sqrt{b(u-3/2)}}{\sqrt{D_0-u}},
 \qquad Y^2=1/2-2/m.
\]
For \(s=\sqrt{u-3/2}<.167\),
\(.3947<b<.3948,\sqrt b>.6282\) give
\[
 (b-.475^2)\sqrt b>.106212915,\qquad
 (b+.475^2)s<.103610975.
\]
Thus \(K_5>.475\).
The payment \(.6/m\leq.12\) fits each
boundary coefficient and gives
\[
 G_{5,5}>1.2(.475-.4208)-.36(.1)=.02904,\qquad
 G_{5,6}>1.2(.475-451/1200)-.36/6=.059.
\]

For native \(7\),
\(.900968<g<.900969\) and \(1.231<u=g^{-2}<1.232\).
The positive actual core
\(c+g^{-1}n_1+2(g^{-1}-g)n_2\), divided by \(g\),
gives
\[
 V=\omega_1-un_1-2(u-1)n_2,\quad
 K_7=\sqrt{D_0-u}-\sqrt{7u-3u^2-4},
\]
\[
 Y^2=6-8u+3u^2-2/m.
\]
Its selected products are zero on \(c,n_1\),
negative on \(n_2,n_3\), and nonpositive on
every unknown wall; time is positive.
On the stated interval,
\(D_0-u>.814^2\) and
\(0<7u-3u^2-4<.070917<.2664^2\), so \(K_7>.545\).
The residual bounds are \(Y_5^2<.299\),
\(Y_6^2<.366\).
Payment \(.4/m\leq.08\) gives gains
\[
 .8(.545-.4208)-.16(.299)=.05152,\qquad
 .8(.545-451/1200)-.16(.366)>.07677.
\]
This proves the table and all its actual-wall
and individual-price hypotheses.
\subsection{The complete two-finite endpoint exclusion}
\label{app:A18finitepair}
Suppose both isolated endpoint owns are finite.
For \(D=36/19,e=2/19\), their actual mutual
magnitude is
\[
 M(g,h)=\sqrt{Dg^2-1}\sqrt{Dh^2-1}-egh.
\]
It increases strictly in each own for
\(g,h\geq g_5=\cos(\pi/5)\).
For example the derivative sign follows by
squaring positive terms from
\(g/\sqrt{Dg^2-1}>1/\sqrt D\) and
\((D^2-e^2)h^2>D\), since
\(D^2-e^2=68/19\) and \(g_5^2>9/17\).

If either own is \(5\), monotonicity gives
\[
 0<\frac{17\sqrt5-25}{76}=M(g_5,g_5)
 \leq M(g_5,h)<M(g_5,1)
 <\sqrt{(9\sqrt5-11)/38}<1/2.
\]
That is forbidden for distinct actual walls.
Thus both own labels \(m,n\) are at least \(7\).
Relabel so \(m\leq n\). Now \(M>0\) already
at \(m=n=7\), while
\(M<M(1,1)=15/19<\cos(\pi/5)\).
The full mutual alphabet therefore leaves
exactly native \(3\) or native \(4\).
Put
\[
 H(g,h)=272g^2h^2-144(g^2+h^2)-8gh+57,
\]
\[
 F_S(g,h)=136g^2h^2-72(g^2+h^2)+19-4Sgh.
\]
They have the signs of \(M-1/2\) and
\(M-1/\sqrt2\), respectively, with
\(S=\sqrt2\). This follows by the
difference of squares against \(egh\)
plus the target; the other factor is
positive. Both increase in each argument
on \([.9,1]^2\), for \(1.4\leq S\leq1.5\),
and \(F_S\) decreases with \(S\).

For mutual \(3\), \(m\geq8\) gives
\[
 M(g_8,g_8)=(17\sqrt2-4)/38>1/2.
\]
At \(m=7\) the exact rational native
enclosures below give
\[
 M(g_7,g_8)<1/2<M(g_7,g_9).
\]
At the upper \(7/8\) pair the bounds
\(u<.8117452,v<.8535543,\ u+v>1.6652993\),
\(uv<.692869,gh>.832387\) yield
\(H<-10018272/10^7<0\).
At the lower \(7/9\) pair,
\(u>.8117433,v>.8830210,\ u+v<1.6947645\),
\(uv>.716786,gh<.846633\) yield
\(H>1146640/10^6>0\).
Thus no integer \(n\geq m\) gives mutual \(3\).

For mutual \(4\), if \(m\leq10\), increase
\(g\) to the upper \(g_{10}^+=.951057\)
and \(h\) to \(1\). Since
\((g_{10}^+)^2<.9045095\) and
\(4(1.414213)g_{10}^+>5.37997\),
\[
 F_{\sqrt2}(g,h)
 \leq F_{1.414213}(g_{10}^+,1)<-.491362<0.
\]
If \(m\geq15\), monotonicity gives
\[
 M\geq(34/19)\cos^2(\pi/15)-1>1/\sqrt2.
\]
Use \(\cos(\pi/15)>10783/11025\),
from \(\cos x>1-x^2/2,\pi<22/7\), and
\(\sqrt2<99/70\). The rational remaining
margin is
\[
 (34/19)(10783/11025)^2-1-99/140
   =1498515515/323324662500>0.
\]
Only \(m=11,12,13,14\) remain.
Their complete consecutive-integer gaps are
\[
\begin{array}{c|c|c}
m&\text{last }n\text{ below target}&
 \text{first }n\text{ above target}\\ \hline
11&31&32\\12&20&21\\13&16&17\\14&15&16.
\end{array}
\]

Here is full rational arithmetic for these gaps.
For \(g=g_m,S=\sqrt2\), set
\[
 A=136g^2-72,\quad B=4Sg,\quad
 C=64g^2-53-B,\quad T=2A-B.
\]
At \(h=1-\delta\),
\(F_S(g,h)=C-T\delta+A\delta^2\).
The native enclosures give
\[
\begin{array}{c|cc|cc}
m&g^2\text{ lower}&g^2\text{ upper}&
 B\text{ lower}&B\text{ upper}\\ \hline
11&.9206248&.9206269&5.42770&5.42773\\
12&.9330110&.9330131&5.46409&5.46412\\
13&.9427264&.9427284&5.49246&5.49249\\
14&.9504826&.9504847&5.51501&5.51504.
\end{array}
\]
Thus the following looser bounds hold:
\[
\begin{array}{c|cc|cc|c}
m&C_-&C_+&T_-&T_+&A_+\\ \hline
11&.49225&.49243&100.981&100.984&53.2053\\
12&1.24858&1.24875&104.314&104.317&54.8898\\
13&1.84199&1.84216&106.928&106.931&56.2111\\
14&2.31584&2.31602&109.016&109.018&57.2660.
\end{array}
\]
For last-below \(n\), use its upper cosine
and \(S=1.414213\).
The \(\delta\) numerators over \(10^6\) are
\(5130,12311,19214,21852\);
upper numerators of \(A\delta^2\) are
\(1401,8400,20800,27374\), also over \(10^6\).
The exact resulting upper numerators over
\(10^9\) are
\[
 -24201530,\ -27059654,\ -191554592,\ -38823632.
\]
For first-above \(n\), use its lower cosine
and \(S=1.414214\).
The \(\delta\) numerators are
\(4816,11170,17027,19215\);
dropping the positive quadratic term gives
lower numerators over \(10^9\)
\[
 5911056,\ 83359110,\ 21275863,\ 221059130.
\]
All signs are strict.
Monotonicity in \(g,h\) proves every integer
\(n\geq m\) lies on one side of the stated
gap, so mutual \(4\) is impossible as well.
This proves the two-finite exclusion.

\subsection{The finite rational angle certificate}
\label{app:A18anglecertificate}
For completeness, the only native enclosures
used in the preceding finite calculation are:
\[
\begin{array}{c|r|r@{\qquad}c|r|r}
j&10^6g_j^-&10^6g_j^+&
j&10^6g_j^-&10^6g_j^+\\ \hline
7&900968&900969&8&923879&923880\\
9&939692&939693&10&951056&951057\\
11&959492&959493&12&965925&965926\\
13&970941&970942&14&974927&974928\\
15&978147&978148&16&980785&980786\\
17&982973&982974&20&987688&987689\\
21&988830&988831&31&994869&994870\\
32&995184&995185&&&
\end{array}
\]
These are exact rational inequalities, not
floating-point evaluations.
Machin's identity
\(\pi=16\arctan(1/5)-4\arctan(1/239)\)
follows by the tangent addition formula
and its first-quadrant sign.
For
\[
 A_s=\sum_{r=0}^s\frac{(-1)^r}{(2r+1)5^{2r+1}},
 \qquad
 B_s=\sum_{r=0}^s\frac{(-1)^r}{(2r+1)239^{2r+1}},
\]
alternating remainders and rational
cross multiplication give
\[
 31415926/10^7<16A_5-4B_2<\pi
 <16A_6-4B_1<31415927/10^7.
\]
For \(0<x<.45\), the alternating cosine
Taylor sums give \(C_{10}(x)<\cos x<C_8(x)\).
Put \(a_j=10^7j,p_-=31415926,p_+=31415927\), and
\[
\begin{split}
Q_{10}(a,p)={}&3628800a^{10}-1814400p^2a^8
 +151200p^4a^6-5040p^6a^4\\
 &+90p^8a^2-p^{10},\qquad Q_8=Q_{10}+p^{10}.
\end{split}
\]
For every displayed row with integers \(L_j,U_j\),
the finite exact certificate is
\[
 10^6Q_{10}(a_j,p_+)-3628800L_ja_j^{10}>0,
\]
\[
 3628800U_ja_j^{10}-10^6Q_8(a_j,p_-)>0.
\]
There are exactly thirty integer inequalities;
their positivity is obtained by direct
integer multiplication of the specified
finite expressions.
The alternating remainder signs and cosine
monotonicity then prove all fifteen rows.
This is a terminating exact verification
recipe with a complete stated output;
no execution of a checker is claimed by
this manuscript transcription.
The historical unused cyclotomic fallback
is not an input.
Finally \(1.414213<\sqrt2<1.414214\)
follows by squaring those two rationals.

\subsection{The complete both-isolated endpoint conclusion}
The local isolated gains, with their separate
prices and disjoint residual supports, exclude
two dotted endpoints:
the three reference-plus-gain lower bounds are
\[
 .9240352+2(.06)>1,\quad
 .8937+.06+4699/60000>1,\quad
 .86414+2(4699/60000)>1.
\]
For one finite and one dotted endpoint, the
full four-pattern totals are
\[
\begin{array}{c|c|c|c|c}
\text{finite own}&5/5&6/5&5/6&6/6\\ \hline
\geq8&1.0311352&>1.01911666&>1.01911666&>1.00787333\\
5&1.0130752&1.0127&>1.00105666&>1.00145666\\
7&1.0355552&>1.03047&>1.02353666&>1.01922666.
\end{array}
\]
In the middle two columns the block order
is finite/dotted. The gains add within
one square because their residuals occupy
opposite zero-sum blocks.
Together with the two-finite theorem,
these exhaust every isolated endpoint own.
Thus both actual endpoint opposites cannot
be isolated.
+
\subsection{The complete actual endpoint inventory}
\label{app:A18endpointinventory}
We now classify an endpoint opposite before using its own contact.
Release \(n_1\). The other endpoint is finite, so the actual Gram on
\(c,n_2,\ldots,n_{18}\) is positive definite. The retained diagram is
the connected \(A_{17}\) path. A one-node finite enlargement is either
disconnected or of type \(A,B,D\): exceptional spherical types have
rank less than eighteen. Thus the new node is isolated, or has one
simple bond to an end or a next-to-end node, or one bond of label four
to an end. Two white contacts would create a cycle and cannot occur
in a finite diagram. This gives the six nonisolated placements below,
without imposing a condition on an arbitrary face.

Write a nonisolated opposite as
\[
 c=A p-g\omega_1-\eta\omega_j,\qquad g=-\ip c{n_1}\geq0,
\]
where \(\eta=1/2\) for a simple white bond and \(\eta=1/\sqrt2\)
for a white bond of label four. Put \(C_{1j}=\ip{\omega_1}{\omega_j}\)
and \(D_j=\omega_j^2\). The actual norm and retained projection give
\[
 A^2=D_0g^2+2g\eta C_{1j}+\eta^2D_j-1,\qquad
 \delta_{\rm white}
 =1-\frac{2\eta^2(j-1)(19-j)}{18}>0.
\]
The edge cosine is
\[
 C=\frac{D_0g+\eta C_{1j}}{\sqrt{D_0\delta_{\rm white}}}.
\]
For a simple end bond the retained projection cost is \(17/36\)
and the residual is \(19/36\); they must not be interchanged.
The next-to-end and label-four residuals are \(1/9\) and \(1/18\).
Substituting in the displayed formula and using
\eqref{app:eq:A18shortends} yields
\[
\begin{array}{c|c|l}
\text{white contact}&C&\text{consequence}\\ \hline
j=2,\ 3&(36g+17)/19&g<103/144:\ \text{own }0,3,4\\
j=3,\ 3&(36g+16)/(2\sqrt{19})&g<.101:\ \text{own }0\\
j=17,\ 3&(36g+2)/(2\sqrt{19})&g<.49:\ \text{own }0\\
j=18,\ 3&(36g+1)/19&g<167/144<7/6\\
j=2,\ 4&(36g+17\sqrt2)/\sqrt{38}&C\geq17/\sqrt{19}>9/4\\
j=18,\ 4&(36g+\sqrt2)/\sqrt{38}&g<.35:\ \text{own }0 .
\end{array}
\]
The numerical cuts follow by positive squaring from
\(\sqrt{19}<4.36,\sqrt{38}<6.165,\sqrt2>1.414\).
A positive own is at least \(1/2\); finite owns are
\(\cos(\pi/m)\), while parallel and dotted owns are respectively
one and greater than one.

For near simple-three with own zero, \(A^2=-2/19\); own three has
the actual positive null sum \(c+n_1+n_2\). For the next-to-near
simple-three with own zero, the actual affine \(E_8\) principal
has center \(n_3\), arms
\[
 (c),\qquad(n_2,n_1),\qquad(n_4,n_5,n_6,n_7,n_8),
\]
and positive kernel coefficients
\[
 6;\qquad3;\qquad4,2;\qquad5,4,3,2,1.
\]
Each row of its unit Gram annihilates this vector. For the
next-to-far simple-three and far label-four with own zero,
\(A^2=-2/19\) and \(-1/19\), respectively.

For far simple-three, own zero and own three have
\(A^2=-10/19\) and zero. Own four has the actual affine path
\[
 (n_{17},n_{18},c,n_1,n_2),\qquad
 \text{bonds }(3,3,4,3),
\]
with positive kernel \((1,2,3,2\sqrt2,\sqrt2)\).
Own six has the null sum \(\sqrt3c+2n_1+n_2\);
parallel has \(c+n_1\). The same nulls exclude isolated
own six and parallel, while isolated owns zero, three and four
have negative time square.

Every positive null combination just used contains its actual
opposite with a positive coefficient. Its time is consequently
positive, so it is a nonzero future null vector. Against any
unselected original wall its product is nonpositive by acuteness;
against its selected walls the kernel equations give zero.
It therefore determines an ideal point of the whole chamber,
contradicting compactness. This checks all original walls,
including those absent from the selected diagram.

The complete surviving list at either end is therefore:
\[
 \begin{array}{ll}
 \text{near:}&\text{simple-three at the neighboring node, own four};\\
 \text{far:}&\text{simple-three at the opposite end, own five,
                finite own at least seven, or dotted};\\
 \text{isolated:}&\text{own five, own seven,
                finite own at least eight, or dotted}.
 \end{array}
\]
In every far case \(g<167/144<7/6\).
The two copies of this classification concern the actual endpoint
walls of the original vertex.

\subsection{Removing the near endpoint and the two-far possibility}
\label{app:A18nearfar}
The remaining near role has
\[
 C=(18\sqrt2+17)/19,\qquad
 x_{\rm end}^2\geq
 \frac{18}{19}\frac{18\sqrt2-2}{18\sqrt2+36}>.601^2.
\]
After clearing positive denominators, the last inequality is
\(200.469258\sqrt2-283.061484>0\); its lower bound at
\(\sqrt2>1.4142\) is \(.4421406636\).
Thus \(x_{\rm end}>.601\).
A strong end gives coordinate gain
\((5/4)(.601-.298)^2=.11476125\);
a weak end gives gain greater than
\((6/5)(.601-.266)^2=.13467\).
Either removes the zero-weak or one-weak reference pattern:
even the smaller one-weak total is
\(.8937+.11476125>1\).

Only two weak ends could remain. A far opposite at the other end
has \(g\geq1/\sqrt2\), a weaker bound than its actual surviving
own list. Its edge inequality gives \(x_{\rm other}>.394\).
The cleared comparison is
\[
 270.909288\sqrt2-382.98968>0;
\]
at \(1.4142\) the margin is \(.1302350896\).
Its weak-block gain is therefore greater than
\((6/5)(.394-.266)^2=.0196608\).
An isolated opposite gives weak gain greater than \(.059\)
by the complete local table; another near opposite gives the
larger near gain. These opposite residuals have disjoint support,
zero cross, and separately paid nonnegative coefficients.
The weakest combined total is
\[
 .86414+.13467+.0196608=1.0184708>1.
\]
Thus no near role survives, on either orientation.

Suppose next that both endpoints are far, and write
\[
 c=A p-g\omega_1-\omega_{18}/2,\qquad
 d=B p-h\omega_{18}-\omega_1/2.
\]
Their time squares are
\[
 A^2=(36g^2+2g-10)/19,\qquad
 B^2=(36h^2+2h-10)/19 .
\]
For \(0<g<7/6\),
\[
 \frac85g^2-A^2
   =\frac{50-10g-28g^2}{95}>\frac2{855}>0,
\]
because the numerator decreases up to \(7/6\).
The same holds for \(h\). Their spatial product satisfies
\[
 \frac{2gh+18g+18h+1/2}{19}
 >\frac{18(g+h)}{19}
 >\frac{216gh}{133}
 >\frac85gh>AB.
\]
Here \(1/g+1/h>12/7\) and \(1080>1064\).
The two caps are distinct: coincidence of their old contact rows
would force \(g=h=1/2\), for which both time squares vanish.
Their positive mutual product contradicts actual acuteness.

The preceding both-isolated theorem and the endpoint inventory
now leave exactly one far endpoint and one isolated endpoint.
Reflect the coordinate formulas if necessary so the far cap is
on the left. An isolated own five is impossible for every far own.
Indeed, with \(u_5=6-2\sqrt5\),
\[
 D_0-u_5=2\sqrt5-78/19<178/475<3/8,\qquad
 A^2<\frac85(7/6)^2=98/45.
\]
Thus \(A^2(D_0-u_5)<49/60<(18/19)^2\).
For an isolated right cap of own \(\sigma\), its time is
\(\sigma\sqrt{D_0-u_5}\), whereas the spatial product divided
by \(\sigma\) is \((18+2g)/19>18/19\).
The actual cap product would be positive. The supports show
these caps are distinct, so this too contradicts acuteness.

\subsection{Uniform removal of the far finite and dotted tail}
\label{app:A18fartail}
A far own at least seven or a dotted own has \(g>.9\).
Its actual length consequently gives
\[
 C>(167/95),\qquad
 x_{\rm far}^2>648/2489>.510^2.
\]
The last positive rational margin is
\(648-2489(.2601)=.6111\).
The legal coordinate gains on the strong and weak blocks are
\[
 G_{\rm far,5}>.05618,\qquad
 G_{\rm far,6}>.0714432.
\]
The other endpoint is isolated with own seven, finite own at least
eight, or dotted. The minimum over the entire local table is
\[
 G_{\rm iso,5}>.0471,\qquad G_{\rm iso,6}>157/2400 .
\]
In particular the weak minimum must not be replaced by the larger
own-seven estimate. The two individually funded residuals belong
to opposite blocks, so their gains add in one square. All four
possibilities give
\[
\begin{array}{c|l}
(m_{\rm far},m_{\rm iso})&
 \text{strict lower bound for }\norm{x}^2\\ \hline
(5,5)&.9240352+.05618+.0471=1.0273152\\
(6,5)&.8937+.0714432+.0471=1.0122432\\
(5,6)&.8937+.05618+157/2400>1.01529666\\
(6,6)&.86414+.0714432+157/2400>1.00099986 .
\end{array}
\]
This removes all these far owns, including arbitrary finite
labels and the continuous dotted interval.

\subsection{The last mutual gap and an actual orthogonal projection}
\label{app:A18jointpoint}
Only far own five remains. Set
\[
 \phi=(1+\sqrt5)/2,\quad g=\phi/2,\quad
 A^2=(10\phi-1)/19,\quad L=(18+\phi)/19,\quad
 \delta=1-3g^2/2=(5-3\phi)/8>0.
\]
Substitution in \(t^2-t-1\) gives \(1.618<\phi<1.619\).
Write the isolated right cap as
\(d=A_Ip-\sigma\omega_{18}\), with \(u=\sigma^{-2}\).
If it were finite, \(u>1\) and
\[
 A^2(D_0-u)<.8(.9)=.72<1<L^2 .
\]
Its mutual product with the far cap would be positive.
Thus \(\sigma>1\) is dotted and \(0<u<1\).
The actual isolated edge has \(C_{\rm iso}=\sqrt{D_0}\sigma<9/4\),
so \(u>64/171\).

The two caps are distinct because the far cap has a positive old
contact at \(n_1\) and the isolated right cap has none.
Their actual mutual product is
\[
 \ip cd=\frac{L-A\sqrt{D_0-u}}{\sqrt u}\leq0.
\]
Since \(A<.9,L>1,D_0-u<260/171<25/16\), and
\(\sqrt u>3/5\), this product is strictly greater than \(-5/24\).
The native alphabet for distinct actual walls has no value in
\((-1/2,0)\). Therefore
\begin{equation}\label{app:eq:A18mutualzero}
 \ip cd=0,\qquad
 AA_I=\sigma L,\qquad
 u=D_0-\frac{L^2}{A^2}
   =\frac{323\phi-361}{19(10\phi-1)}.
\end{equation}
Orthogonality is proved here before it is used.
The rational function has derivative with positive numerator
\(62453\). At \(1.618,1.619\) its values are
\(80807/144210\) and \(161937/288610\).
The positive comparisons
\[
 25(80807)-14(144210)=1235,\qquad
 281(288610)-500(161937)=130910
\]
give
\begin{equation}\label{app:eq:A18uinterval}
 .560<u<.562 .
\end{equation}
These are geometric identities; no integrality or field restriction
is being inferred from them.

Consider the positive actual correction
\[
 \eta=c+g(3n_1/2+n_2+n_3/2)
      =A p-g\omega_4/4-\omega_{18}/2 .
\]
Direct products give
\[
 \ip\eta{n_1}=\ip\eta{n_2}=\ip\eta{n_3}=0,\qquad
 \eta^2=\ip\eta c=\delta>0.
\]
Thus \(\eta\) is spacelike. The actual \(H_4\) block
\(c,n_1,n_2,n_3\) is positive definite, its Schur complement being
\(\delta\). Equation \eqref{app:eq:A18mutualzero} and the right
cap's isolation make \(d\) orthogonal to this block.

The required chamber point is the projection
\[
 q_0=p+(A/\delta)\eta+A_Id,\qquad
 H=A^2/\delta+A_I^2,\qquad
 q_0^2=\ip p{q_0}=-1-H .
\]
It has zero product with every selected wall. Every correction
coefficient on an actual normal is positive; against every
unselected original wall, acuteness and \(\ip p a\leq0\)
therefore give \(\ip{q_0}a\leq0\). This includes \(n_{18}\)
and all unknown facets. Also \(\ip p{q_0}<0\), so \(q_0\) is
future timelike. Its normalization is a point of the whole
chamber. The spacelike vector \(\eta\) itself was never used
as a chamber point.

Put
\[
 a=\phi^{-1}+\frac{\delta L}{A^2gu}>0,\qquad
 V=\omega_4/4+a\omega_{18},\qquad
 U=A^2+\delta D_0/u .
\]
Using \(AA_I=\sigma L\), the point has the full decomposition
\[
 q_0=(1+H)p-(Ag/\delta)V,\qquad
 \delta(1+H)=U.
\]
Compare its normalization with the same original farthest vertex:
\[
 -\ip z{q_0/\sqrt{1+H}}\leq-\ip zp=E .
\]
Substitution of \(z=E(p-x)\) gives the actual pressure
\begin{equation}\label{app:eq:A18jointpressure}
 \ip xV\geq K,\qquad
 K=\frac{\delta(1+H-\sqrt{1+H})}{Ag}
   =\frac{U-\sqrt{\delta U}}{Ag}.
\end{equation}
No auxiliary point has been declared farthest.

All subsequent estimates have positive denominators. The enclosure
of \(\phi\) and \eqref{app:eq:A18uinterval} give
\[
\begin{gathered}
 .7989<A^2<.7995,\qquad .0178<\delta<.0183,\\
 .809<g<.8095,\quad A<.895,\quad L<1.033 .
\end{gathered}
\]
For the lower estimate retain \(D_0=36/19\) exactly:
\[
 \frac{\delta D_0}{u}
  >\frac{.0178(36/19)}{.562}
   =\frac{3204}{53390}>.06 ,
\]
where \(50(3204)-3(53390)=30>0\).
The upper bound is \(\delta D_0/u<.0625\), by positive
cross multiplication. Consequently
\[
 .8589<U<.862,\quad
 \sqrt{\delta U}<.126,\quad
 Ag<.7245025,\quad
 K>\frac{.7329}{.7245025}>1.01 .
\]
For the radical comparison,
\(.0183(.862)=.0157746<.126^2\), and
\(1.01(.7245025)=.731747525<.7329\).
Finally,
\[
 \delta L<.0189039
 <.06(.7989)(.809)(.560)=.02171601936 .
\]
Since \(\phi^{-1}<.619\), this proves \(0<a<.679<.68\).

\subsection{One final funded vector for every end pattern}
\label{app:A18lastpayment}
For the pressure \eqref{app:eq:A18jointpressure}, subtract
\[
 B=\omega_{m_L}/m_L+a\omega_{19-m_R}/m_R,\qquad Y=V-B.
\]
The left residual is \(\omega_4/4-\omega_{m_L}/m_L\);
the right is \(a\) times the reflected
\(\omega_1-\omega_{m_R}/m_R\).
They are supported on the first \(m_L\) and last \(m_R\)
coordinates, have zero sum on each block, and have zero cross.
Their full squared norm is exactly
\[
 Y^2=\frac12-\frac2{m_L}
          +a^2\left(2-\frac2{m_R}\right).
\]
In particular the nonzero cross between the original endpoint
duals has disappeared only after the proper block subtractions.

Use the single nonnegative weight \(\mu=.4\).
It costs \(.4/m_L\leq.08\) at the left anchor and
\(.4a/m_R<.0544\) at the right, each below its own available
coefficient \(.1349\) or \(>.1339\).
The all-hole reference and the individual-price hypotheses of
the completed-square rule are therefore satisfied simultaneously.
Write
\(\ell_5=H_5/5=.4208\) and
\(\ell_6=H_6/6=451/1200\).
The gain is strictly greater than
\[
 .8(1.01-\ell_{m_L}-.68\ell_{m_R})
 -.16\left[\frac12-\frac2{m_L}
             +.68^2\left(2-\frac2{m_R}\right)\right].
\]
Replacing \(a\) by its upper bound lowers the linear deficit
and increases the norm penalty. The four rational budgets are
\[
\begin{array}{c|c|c}
(m_L,m_R)&\text{gain lower bound}&
 \text{strict lower bound for }\norm{x}^2\\ \hline
(5,5)&.1080704&1.0321056\\
(6,5)&>.13337706&>1.02707706\\
(5,6)&.1276&>1.0213\\
(6,6)&>.15290666&>1.01704666 .
\end{array}
\]
For example the last bracket equals \(703/750\) and its
linear deficit is \(.3786\); the unrounded gain is
\(.8(.3786)-.16(703/750)\).
Every pattern contradicts \(\norm{x}^2<1\).
One original vector, one reference and one set of individually
available coefficients were used throughout.

\begin{proposition}\label{app:prop:wholeA18}
An actual original farthest vertex of a compact hyperbolic Coxeter
polytope in dimension eighteen cannot have entire incident diagram
\(A_{18}\).
\end{proposition}
\begin{proof}
The finite-white inventory leaves only near, far or isolated
endpoint opposites. The near role is impossible on either side.
Both isolated endpoints are excluded by their complete local and
finite-pair arguments; both far endpoints violate acuteness.
Thus precisely one endpoint is far and the other isolated.
The uniform coordinate/local-gain argument removes every far
own at least seven and every dotted far own. In the remaining
far-own-five case, every finite isolated opposite violates
acuteness. The dotted isolated opposite has mutual product zero
by the actual native gap, and the actual \(H_4+A_1\) point gives
the final simultaneous-payment contradiction above.
These alternatives exhaust the actual endpoint walls, without
arithmeticity, a finite-label cutoff, arbitrary-face heredity,
or a rank decrement.
\end{proof}





